\section{Normalized traces and quantifier-free loop reflection}
\label{sec:traces}

Many mathematical proofs are finite algorithms with an invariant: Euclid's
algorithm, Gaussian elimination under pivots, recurrence evaluation, finite
search, interval refinement, dynamic programming, and certificate checking.
The obstacle in \(\iPRA\) is not iteration itself---primitive recursion is
available---but keeping the induction matrix quantifier-free.  DARC therefore
normalizes each bounded loop to explicit state registers and reflected bits.

\subsection{Trace normal form}

A normalized loop block consists of:
\begin{enumerate}
  \item a finite tuple of state terms
        \(s=(s_1,\ldots,s_k)\);
  \item concrete primitive-recursive initial values \(s^0_j(\vec x)\);
  \item concrete primitive-recursive step terms
        \(T_j(s_1,\ldots,s_k,\vec x)\);
  \item a quantifier-free invariant \(I(s,\vec x)\);
  \item an optional quantifier-free exit/output relation.
\end{enumerate}
The trace functions \(S_j(t,\vec x)\) are generated by simultaneous primitive
recursion.  If the library's simultaneous recursion is binary, larger tuples
are encoded by checked pairing and projections.  Their defining equations are
proved at arbitrary arguments:
\begin{align*}
  S_j(0,\vec x)&=s^0_j(\vec x),\\
  S_j(S t,\vec x)&=T_j(S_1(t,\vec x),\ldots,S_k(t,\vec x),\vec x).
\end{align*}

\subsection{Removing bounded formulas from invariants}

Suppose an informal invariant contains
\(\forall i<n\,P(i,s)\).  Before loop reflection, compile this bounded
formula to a bit \(b_P(n,s)\) by \cref{thm:bounded-reflection} and replace the
invariant clause by \(b_P(n,s)=1\).  Likewise, a bounded existential is
replaced by its finite disjunction bit, and a finite data-structure validity
condition is replaced by a reflected checker bit.  The resulting formula is
syntactically quantifier-free.

This transformation is stronger than observing that bounded quantifiers are
``morally finite.''  The induction axiom is syntax-sensitive; only the
explicit bit equality qualifies.

\subsection{Trace invariant theorem}

\begin{theorem}[Normalized trace reflection]
\label{thm:trace-reflection}
Let \(S(t,\vec x)\) be the generated state tuple and let \(I\) be
quantifier-free.  If ordinary derivations establish
\[
  \vdash I(S(0,\vec x),\vec x)
\]
and
\[
  I(s,\vec x)\vdash I(T(s,\vec x),\vec x),
\]
then the ordinary theory proves the open invariant
\[
  \vdash I(S(t,\vec x),\vec x).
\]
\end{theorem}

\begin{proof}
Substitute the trace terms into the invariant and use their zero equations to
obtain the base case.  In the step case, assume
\(I(S(t,\vec x),\vec x)\), instantiate the local preservation theorem at the
current trace, and rewrite the result by the successor equations to obtain
\(I(S(S t,\vec x),\vec x)\).  Apply quantifier-free induction
\cref{lem:qfind}.  Quantifier-freeness is preserved under term substitution.
\end{proof}

\subsection{Trace certificates versus trace functions}

An external proof may arrive with a list of all states.  DARC distinguishes two
uses.

\paragraph{Recomputation.}
When the step is deterministic and cheap, define the trace by primitive
recursion and prove the invariant as above.  No trace certificate needs to
escape.

\paragraph{Checked trace.}
When a certificate stores auxiliary choices, encode a bounded sequence
\(c\) and define a primitive-recursive checker for the initial state,
transitions, and final condition.  Reflection proves that an accepted trace
has the required property.  If the public theorem does not mention the trace,
seal it existentially.  If it asks for replay data, activate the trace
certificate coordinate and export a bounded canonical code.

The deterministic and certified modes may coexist: deterministic registers
are recomputed while noncanonical hints remain in the local certificate.

\subsection{Demand-directed state slicing}

A loop state can contain many registers.  Backward demand determines which
registers and precision levels affect the requested output.  Let
\(D\subseteq\{1,\ldots,k\}\) be the demanded registers.  A sound slice must
satisfy closure: if \(T_j\) for \(j\in D\) reads register \(i\), then
\(i\in D\).  Iterating this register dependency is another finite monotone
closure and can be folded into the global resource lattice.

\begin{proposition}[Loop slicing]
\label{prop:loop-slicing}
If demanded registers are closed under the read dependencies of their step
terms and the sliced invariant mentions only those registers, then the trace
reflection theorem can be instantiated on the sliced state without retaining
undemanded registers.
\end{proposition}

\begin{proof}
The closed read condition makes every demanded next-state term definable from
the sliced current state.  Generate simultaneous recursion only for the slice.
The base and step proofs for the sliced invariant are the corresponding
specializations of the original local contracts.  Apply
\cref{thm:trace-reflection}.
\end{proof}

The proposition is conditional on having local proofs that do not secretly use
an erased register.  The transfer summary must record all such dependencies.

\subsection{Termination and fuel}

DARC never derives termination of an unbounded partial algorithm merely from a
trace checker.  A loop node has one of three forms:
\begin{enumerate}
  \item fixed or input-bounded fuel, so the trace exists by primitive
        recursion;
  \item a proved primitive-recursive ranking bound, such as the existing
        Euclidean-algorithm bound;
  \item an explicit termination-bound residual.
\end{enumerate}
The final theorem reflects exactly the available form.  A fuel-bounded
statement may be mathematically useful even when a uniform termination theorem
is unavailable.

\subsection{Worked invariant: sum of odd numbers}

Let the state be \((n,a)\), initialized by \((0,0)\), with step
\[
  (n,a)\longmapsto(Sn,a+2n+1).
\]
The invariant is the polynomial equation \(a=n^2\).  At the step,
\[
  a+2n+1=n^2+2n+1=(n+1)^2.
\]
This identity is handled by the checked semiring/ring infrastructure, and the
invariant is quantifier-free.  The reference example executes five steps and
checks the states; a full object proof is a direct instance of
\cref{thm:trace-reflection}.

This simple example illustrates the intended division of labor: Python can
show the generated fixture, but the mathematical theorem is the ordinary
induction derivation.
